\documentclass[../../main.tex]{subfiles}
\begin{document}

\begin{center}
  {\fontsize{48pt}{58pt}\selectfont \textbf{1972}\normalsize\textbf{年度}}
\end{center}

\begin{table}[htb]
  \centering
  \caption{出題テーマと難易度}
  \begin{tabular}{|c|p{6.2cm}|c|c|c|}
    \hline
    \textbf{問題番号} & \textbf{テーマ} & \textbf{思考量} & \textbf{計算量} & \textbf{総合} \\
    \hline
    第1問 & 空間内の二点を見込む角と平面上の面積 & 普 & 普 & 普 \\
    \hline
    第2問 & 鏡映したベクトルと三角形の角 & 普 & 普 & 普 \\
    \hline
    第3問 & 接する円の列の半径の極限 & 普 & 普 & 普 \\
    \hline
    第4問 & 三次方程式の実根の個数 & 普 & 普 & 普 \\
    \hline
    第5問 & 積分で定まる関数と二階微分 & 普 & 普 & 普 \\
    \hline
    第6問 & 道路上の確率と警官の最適配置 & 普 & 難 & 普 \\
    \hline
  \end{tabular}
\end{table}

\begin{table}[htb]
  \centering
  \caption{解答}
  \renewcommand{\arraystretch}{1.4}
  \begin{tabular}{|c|c|p{9cm}|}
    \hline
    \textbf{問題番号} & \textbf{小問} & \textbf{答え} \\
    \hline
    第1問 & -- & 面積は$13\pi+75\sqrt{3}$ \\
    \hline
    第2問 & -- & $(m,\angle BAC,|\mathbf{c}|)=\left(1,\dfrac{2\pi}{3},1\right),\left(2,\dfrac{3\pi}{4},\sqrt{2}\right),\left(3,\dfrac{5\pi}{6},\sqrt{3}\right)$ \\
    \hline
    第3問 & -- & 極限値は$R$ \\
    \hline
    第4問 & -- & $k=1$なら$0$個，$k=0,4$なら$2$個，$0<k<4$かつ$k\ne1$なら$1$個，$k<0$または$k>4$なら$3$個 \\
    \hline
    第5問 & -- & $h(x)=-x$ \\
    \hline
    第6問 & -- & $P_4,P_6,P_7$に配置．最大確率は$\dfrac{17}{32}\fallingdotseq0.53$ \\
    \hline
  \end{tabular}
\end{table}

\paragraph{試験形態}
6問．試験時間は要確認．

\end{document}
